\chapter{Apéndice B: Uso ético de herramientas de IA}
\label{Appendix:2}

Este apéndice complementa la ``Declaración responsable sobre autoría y uso ético de herramientas de IA'' proporcionada por la UCM y que se adjunta con esta memoria. En cumplimiento de los principios de transparencia y ética académica, además del cumplimiento de las normativas europeas sobre el uso de inteligencia artificial en los textos de interés público \citep{CodeOfPracticeAI2026}. Declaramos aquí el uso de herramientas y algoritmos basados en inteligencia artificial durante el desarrollo de este trabajo:

\begin{itemize}
	\item \textbf{Accesibilidad y transcripción:} Debido a un problema medico personal del autor, diagnosticado durante el desarrollo del proyecto, la cual limitaba severamente la capacidad de escritura manual prolongada mediante teclado, se emplearon herramientas de reconocimiento de voz a texto (concretamente, el servicio de dictado nativo de Windows, atajo \textit{Windows + H}) para la redacción de los borradores de esta memoria o texto completo que en todo momento fue supervisado.
	
	\item \textbf{Revisión y corrección de estilo:} Se utilizaron herramientas de asistencia a la redacción para corregir errores tipográficos y depurar la gramática del texto, asi como la revisión del texto a través de LLms para revisar huecos o frases incompletas que se hubieran pasado por alto, sin que estas herramientas alteraran en ningún momento la originalidad de las ideas o el contenido intelectual del autor.
	
	\item \textbf{Modelado analítico y tratamiento de datos:} Tal y como se detalla en el \hyperref[cap:Materiales y métodos]{Capítulo \ref*{cap:Materiales y metodos}}, el núcleo analítico del trabajo se basa en la aplicación de algoritmos de aprendizaje automático (\textit{Machine Learning}). Se han utilizado herramientas de IA orientadas al análisis de datos (como \textit{K-Means} o agrupamiento jerárquico) para extraer conclusiones, identificar patrones y generar los resultados de los casos de estudio.
	
	\item \textbf{Asistente de GCP:} \textit{Google Cloud} cuenta con un asistente de inteligencia artificial integrado en la consola para el análisis de logs (\textit{Cloud Assist}), el cual se utilizó para intentar solucionar las incidencias durante los despliegues iniciales en \textit{Cloud Run} y \textit{Cloud Functions}. Aunque aportó pistas sobre el comportamiento de los contenedores, no se termino de encontrar una solucion viable y el trabajo se orientó hacia el uso de \textit{Vertex AI Workbench} para unificar de forma fluida el procesamiento y el análisis.
	
	\item \textbf{Modelos de lenguaje y motores de búsqueda avanzados:} Se emplearon de forma complementaria modelos de lenguaje de grandes dimensiones (\textit{LLMs}) integrados en buscadores y plataformas oficiales de documentación técnica, como una evolución avanzada de foros de consultas tipo \textit{Stack Overflow}. Su uso se limitó a la resolución de incidencias puntuales, validación de parámetros de las CLI de AWS y GCP, y la solución puntual de errores durante las fases de despliegue.
	
	\item \textbf{:} %[TODO] seguro que uso mas, vertex y tal se incluye? 
\end{itemize}

\newpage
\section{Análisis de la gestión de datos mediante LLMs}
\label{sec:Apendice_LLMs}

% [TODO] desarrollar estudio adicional sobre cómo los LLMs pueden manejar los datos??????

